\section{Related Work}
\label{sec:related}

\para{Similarity indices for neural representations.} \citet{kornblith2019similarity} introduced \ckam (Centered Kernel Alignment) as a similarity index that is invariant to orthogonal rotation and isotropic scaling but not to arbitrary invertible linear maps. That choice of invariance is forced by Theorem 1 of the same paper, which shows that any index invariant under invertible linear transformations becomes degenerate when feature width exceeds dataset size---ruling out CCA for wide modern networks~\citep{raghu2017svcca}. Subsequent work proposed alternative indices: Normalized Space Alignment~\citep{ebadulla2024nsa}, graph- and topology-based measures~\citep{hryniowski2020nnts}, and sample-wise cosine similarity~\citep{jiang2024tracing}. \citet{davari2022reliability} showed that \ckam can be manipulated by translations of small subsets without changing the underlying function, motivating multi-metric confirmation. We follow this line: we report \ckam alongside orthogonal Procrustes and sample-cosine, and show that the three metrics rank layer pairs consistently ($\rho = 0.98$ for Procrustes, $0.72$ for cosine).

\para{Cross-layer similarity in CNNs, ResNets, and transformers.} \citet{nguyen2021wide} coined the ``block structure'' phenomenon: in over-parameterized ResNets, large contiguous runs of layers have \ckam close to 1. \citet{nguyen2022origins} traced this to a small set of ``dominant datapoints'' with high activation norm, showing that excluding the top-decile reduces block-structure \ckam by $\Delta \approx\!-0.15$ in CNNs. \citet{raghu2021vit} reported that vision transformers show much more uniform cross-layer CKA than CNNs, driven by skip connections. \citet{jiang2024tracing} proved that in residual networks sample-wise cosine similarity tracks \ckam, and that the final classifier can be applied to any hidden layer---the logit-lens. \citet{brown2023dissimilarity} extended block structure to generative transformers. \citet{pessoa2023onewide} showed that transformer FFNs across layers are highly redundant and can be shared. Our work maps these same phenomena to pure MLPs on image benchmarks, finding a weaker dominant-datapoint effect ($\Delta = -0.011$) but clear block structure at depth 16.

\para{MLP-specific analyses.} Pure MLPs are under-tabulated in the representation-similarity literature. \citet{hryniowski2020nnts} studied LeNet-5 on \mnist using k-NN topological similarity. \citet{ding2025selfsimilarity} introduced self-similarity as a graph-based metric and found that MLPs, CNNs, and attention models differ in their self-similarity profiles, with an $SS_{\mathrm{rate}}$ regularizer improving MLP accuracy by up to 6\%. \citet{tolstikhin2021mlpmixer} re-popularized deep MLP stacks with MLP-Mixer, whose patch-mixing FFNs invite exactly the cross-layer similarity questions we pose. Neither of these works produces a full depth$\times$width heatmap grid for classical MLPs on \cifar.

\para{Rich vs.\ lazy feature learning.} \citet{chizat2019global} quantifies the transition between the lazy (NTK-like) regime, where trained weights stay close to initialization, and the rich regime, where representations deviate substantially. \citet{tong2025richness} predicted that lazy MLPs keep cross-layer representations close to random initialization---trivially similar---while rich MLPs learn distinct per-layer features. We empirically confirm this prediction by sweeping the initialization scale and measuring $\|\theta_{\mathrm{final}} - \theta_{\mathrm{init}}\|/\|\theta_{\mathrm{init}}\|$ against mean off-diagonal CKA: lazy runs reach CKA $= 0.76$ but only 40\% test accuracy.

\para{Practical consequences of cross-layer redundancy.} MPruner~\citep{park2024mpruner} uses inter-layer CKA to prune redundant layers; \citet{pessoa2023onewide} share a single FFN across transformer layers with $<$1 BLEU loss; the LAYA line of work aggregates predictions across intermediate layers. Our block-structure finding for deep MLPs suggests the same pruning opportunities in classical MLPs, which we flag as follow-up work.
